\section{Modelo matemático del PMSM}

\subsection{Dinámica eléctrica en $dq$}

La representación de referencia empleada para la planta es:
\begin{align}
\frac{di_d}{dt} &= \frac{v_d-R_s i_d+\omega_e L_q i_q}{L_d},\\
\frac{di_q}{dt} &= \frac{v_q-R_s i_q-\omega_e(L_d i_d+\lambda_m)}{L_q}.
\end{align}

\subsection{Par electromagnético}

Como formulación clásica de referencia:
\begin{equation}
T_e=\frac{3}{2}p\left(\lambda_m i_q+(L_d-L_q)i_d i_q\right).
\end{equation}

Esta expresión se trata como ecuación de referencia para la reconstrucción y no como identificación exacta de la ley interna del modelo Simscape histórico.

En la parametrización actual $L_d=L_q$, por lo que el término clásico de reluctancia se anula.

\subsection{Dinámica mecánica}
\begin{align}
\frac{d\omega_m}{dt} &= \frac{T_e-T_{load}}{J},\\
\frac{d\theta_m}{dt} &= \omega_m.
\end{align}

La relación entre ángulo mecánico y eléctrico es:
\begin{equation}
\theta_e=p\theta_m=4\theta_m.
\end{equation}
